\chapter{Grundlagen}
\label{ch:grundlagen}

\section{Digitale Bildverarbeitung}
\label{sec:bildverarbeitung}
% TODO: Bildaufnahme, Graustufen, Filterung, Schwellwertverfahren nach Otsu~\cite{otsu1979},
% Morphologie, Konturen, Bildmomente. Verwendete Bibliothek: OpenCV~\cite{bradski2000opencv}.

\section{Kameramodell und Kalibrierung}
\label{sec:kalibrierung}
% TODO: Lochkameramodell, intrinsische/extrinsische Parameter, Linsenverzeichnung,
% Kalibrierung nach Zhang~\cite{zhang2000}. Vertiefung in~\cite{hartley2004}.

\section{Fiducial Marker und Homographie}
\label{sec:aruco}
% TODO: ArUco-Marker~\cite{garrido2014}, Detektion, Homographie zwischen zwei Ebenen
% (Direct Linear Transform)~\cite{hartley2004}, perspektivische Entzerrung.

\section{Koordinatentransformationen}
\label{sec:transformationen}
% TODO: Homogene Koordinaten, Starrkörpertransformation, Bestimmung aus Punktkorrespondenzen
% nach der Methode der kleinsten Quadrate~\cite{umeyama1991}.

\section{Industrierobotik und Greifen}
\label{sec:robotik}
% TODO: Kinematik, Koordinatensysteme (Basis, Werkzeug/TCP), Bewegungsarten PTP/LIN,
% Greifertypen, Greifplanung~\cite{siciliano2016}. Kollaborative Roboter und Sicherheit.

\section{Optische Messtechnik und Prüfmerkmale}
\label{sec:messtechnik}
% TODO: Beleuchtung, Kantendetektion, Kreisdetektion mit Hough-Transformation~\cite{duda1972},
% Kalibrierung px -> mm, Messunsicherheit.

\section{Texterkennung (OCR)}
\label{sec:ocr}
% TODO: Ablauf einer OCR-Pipeline, Tesseract~\cite{smith2007}, lernbasierte Verfahren.

\section{Lernbasiertes Greifen}
\label{sec:ki-greifen-grundlagen}
% TODO (Person 4): Überblick lernbasierter Greifansätze, z.\,B. GG-CNN~\cite{morrison2018}.
